\subsection{GPT-2 的 WebText 和 OpenWebText}

GPT-2 的 WebText 思路很代表性：从 Reddit 中高 karma 帖子的外链网页出发，用”人类实际愿意转发什么”作为粗糙质量代理。具体来说，WebText 收集了 Reddit 上获得 3 个以上 karma 的帖子中的外链网页，共 800 万页、约 40GB 文本。开源复现版本 OpenWebText 再通过 URL 抽取、英文过滤（使用 Facebook 的 fastText 分类器）和近重复删除来构建类似数据集。

WebText 的处理流水线可以总结为四步：
\begin{enumerate}
    \item \textbf{社交过滤}：只保留 Reddit karma $\geq 3$ 的帖子中的外链 URL。
    \item \textbf{URL 抽取}：从帖子元数据中提取所有外链网页地址。
    \item \textbf{语言过滤}：用 fastText 分类器过滤掉非英文页面。
    \item \textbf{近重复删除}：移除内容高度相似的页面，避免冗余。
\end{enumerate}

\begin{importantbox}{WebText 的方法论意义}
它不是在证明 Reddit 等于高质量，而是在证明一个更现实的点：\textbf{当全网太大、太乱时，人类行为信号可以作为粗糙的质量代理。}
\end{importantbox}

